% GL Experiments — WebGL learning presentation
% Build (from this folder): pdflatex presentation.tex
% From repo root: pdflatex presentation/presentation.tex
% (run pdflatex twice if TOC page numbers look wrong)

\documentclass[aspectratio=169]{beamer}

\usetheme{Madrid}
\usecolortheme{default}

\usepackage[utf8]{inputenc}
\usepackage[T1]{fontenc}
\usepackage{hyperref}
\usepackage{listings}
\usepackage{xcolor}

\lstdefinestyle{glsl}{
  basicstyle=\ttfamily\footnotesize,
  keywordstyle=\color{blue!70!black},
  commentstyle=\color{gray},
  stringstyle=\color{teal},
  breaklines=true,
  showstringspaces=false,
  frame=single,
  framerule=0.2pt,
  columns=fullflexible,
}
\lstset{style=glsl}

\title{GL Experiments}
\subtitle{WebGL shaders, 2D transforms, and the raster pipeline}
\author{Jacob}
\date{\today}

\begin{document}

\begin{frame}
  \titlepage
\end{frame}

\begin{frame}{Outline}
  \tableofcontents
\end{frame}

\section{Why WebGL}

\begin{frame}{What this project is}
  \begin{itemize}
    \item Hands-on repo for \textbf{OpenGL / WebGL} basics: buffers, shaders, 2D math.
    \item Two tracks: \textbf{browser WebGL} (\texttt{webgl/}) and a small \textbf{C++ OpenGL} app (WIP).
    \item Goal: connect \textbf{pixel-space geometry} to \textbf{GPU drawing} with minimal magic.
  \end{itemize}
  \vfill
  \begin{block}{Primary tutorial series}
    \href{https://webglfundamentals.org/}{WebGL Fundamentals} — translation, rotation, scale, matrices, then 3D.
  \end{block}
\end{frame}

\section{The GPU pipeline}

\begin{frame}{Vertex shader vs.\ fragment shader}
  \begin{description}
    \item[Vertex shader] Runs once \textbf{per vertex}. Must write \texttt{gl\_Position} (clip space). Can output \texttt{varying}s to interpolate.
    \item[Fragment shader] Runs once \textbf{per pixel} (fragment) inside drawn triangles. Writes color (\texttt{gl\_FragColor} in WebGL1 GLSL).
  \end{description}
  \vfill
  \centering
  \small
  CPU: buffer data $\rightarrow$ \texttt{drawArrays} $\rightarrow$ GPU: vertices $\rightarrow$ rasterize $\rightarrow$ fragments
\end{frame}

\begin{frame}{Attributes and uniforms}
  \begin{description}
    \item[Attribute] Per-vertex input (e.g.\ \texttt{a\_position} from a VBO). Location bound with \texttt{vertexAttribPointer}.
    \item[Uniform] Same value for the whole draw call (e.g.\ resolution, translation, rotation).
  \end{description}
  \vfill
  \begin{block}{In this repo}
    The ``F'' shape is a single \texttt{Float32Array} of 2D positions; sliders update uniforms and \texttt{drawScene()} redraws.
  \end{block}
\end{frame}

\section{Our 2D vertex shader}

\begin{frame}{Transform order (concept)}
  We compose a 2D rigid transform using \textbf{homogeneous coordinates} (\texttt{mat3}, \texttt{vec3(x,y,1)}):
  \[
    T_{\text{trans}} \cdot T_{\text{center}} \cdot R \cdot T_{-\text{center}} \cdot S
  \]
  \begin{itemize}
    \item \textbf{S} — scale around local origin
    \item \textbf{R} — rotate around \texttt{u\_rotationCenter}
    \item \textbf{T} — translate in pixel space
  \end{itemize}
  \vfill
  \small
  Same building blocks as \href{https://webglfundamentals.org/webgl/lessons/webgl-2d-matrices.html}{WebGL Fundamentals: 2D matrices}.
\end{frame}

\begin{frame}[fragile]{Vertex shader (excerpt): matrices}
\begin{lstlisting}[language=C]
// u_rotation = vec2(sin(theta), cos(theta))
mat3 rotationMatrix = mat3(
  u_rotation.y, u_rotation.x, 0.0,
 -u_rotation.x, u_rotation.y, 0.0,
  0.0,          0.0,          1.0
);
vec3 worldPosition =
  translationMatrix * toCenter * rotationMatrix
  * fromCenter * scaleMatrix * vec3(a_position, 1.0);
\end{lstlisting}
\end{frame}

\begin{frame}[fragile]{Pixel space to clip space}
  After world position in \textbf{pixels} (origin top-left style), map to clip space $[-1,1]$:
\begin{lstlisting}[language=C]
vec2 zeroToOne = position / u_resolution;
vec2 zeroToTwo = zeroToOne * 2.0;
vec2 clipSpace = zeroToTwo - 1.0;
gl_Position = vec4(clipSpace * vec2(1.0, -1.0), 0.0, 1.0);
\end{lstlisting}
  \texttt{vec2(1,-1)} flips Y: canvas Y grows downward; clip space Y grows upward.
  \vfill
  \small
  See the resolution / clip-space discussion in the WebGL Fundamentals 2D articles.
\end{frame}

\begin{frame}[fragile]{Fragment shader (current)}
  Solid color for every pixel inside the triangles:
\begin{lstlisting}[language=C]
precision mediump float;
void main() {
  gl_FragColor = vec4(1.0, 0.0, 0.5, 1.0);
}
\end{lstlisting}
  Next step in many tutorials: pass \texttt{varying} colors or texture coordinates from the vertex shader.
\end{frame}

\section{Running the demo}

\begin{frame}{Local WebGL demo}
  \begin{enumerate}
    \item \texttt{cd webgl}
    \item \texttt{python3 -m http.server 8080}
    \item Open \texttt{http://localhost:8080/}
  \end{enumerate}
  \vfill
  \begin{itemize}
    \item Swap \texttt{vertex-*.js} imports in \texttt{main.js} to try translation-only, rotation, clip-space variants, etc.
    \item Sliders: translation, scale, rotation (unit circle UI).
  \end{itemize}
\end{frame}

\section{References}

\begin{frame}{Useful links}
  \begin{itemize}
    \item \href{https://webglfundamentals.org/}{WebGL Fundamentals} — main tutorial path.
    \item \href{https://learnopengl.com/}{Learn OpenGL} — often used for desktop GL; concepts overlap.
    \item Notion notes (project): add your workspace link here if you want it in the PDF.
  \end{itemize}
\end{frame}

\begin{frame}{}
  \centering
  \Large Questions?
\end{frame}

\end{document}
